\subsection*{Every fixed prime-power threshold}
The square cutoff is the first of a complete sequence. For an optimizer with primes compressed and exponents nonincreasing, let $w_j$ be the largest prime whose exponent is at least $j$ (zero if there is none). For each fixed positive integer $j$, put $\beta_j=2/(j+1)$. Then, as $k\to\infty$,
\[
                   w_j\sim \beta_j^{1/(j+1)} q^{\beta_j}.
\]
The conclusion is uniform in the choice among optimizing compressed periods for each fixed $j$.

\begin{proof}
The cases $j=1,2$ were proved above. Write $t_j=\beta_j^{1/(j+1)}q^{\beta_j}$ and $\lambda_q=\log q/q^2$. For a period $L$ set
\[
 W_x(L)=\sum_{1<d\mid L}(1/d-1/x)_+.
\]
If $T_L$ is its $k$th nonunit divisor, then $S_L=k/T_L+W_{T_L}(L)$ and $S_L\le k/x+W_x(L)$ at every $x>0$. A change that strictly increases $W_x$ for every $k\le x\le Ck$ improves $S_L$ provided the changed period still has $T_L$ in this range.

For an optimizer the total reciprocal mass of nondivisors among the integers $1\le n\le q$ is $o(1)$. Indeed all these are smooth and the inequality for the truncated missing mass in the proof of Theorem 2 bounds their harmonic mass by $(V(k)-U(k))/(1-q/T)=o(1)$. All the changes below alter divisibility only at powers larger than $q$. They preserve this fact. They also preserve all first-power primes. For the original and changed lattices, only $o(k)$ of the $q$-smooth numbers up to $Ck$ can fail to divide the period for each fixed $C$. The absent primes through $q$ have total reciprocal mass $o(1)$. For each available prime $p$ with exponent $e_p\ge1$, the first obstructing power has reciprocal at most $p^{-2}$. Its contribution tends to zero at every fixed prime since every fixed prime power eventually divides the period. Dominated convergence makes the sum of these reciprocals $o(1)$, so a union bound on multiples proves the claim. Choosing $C>1/\rho$ now gives $k\le T_L\le Ck$ for the original and changed lattices.

Consider changing the inclusion of $p^i$, with $i\ge2$ fixed and $p/t_i$ in a fixed compact subinterval of $(0,\infty)$. Removing this power means decreasing the exponent to $i-1$, even if it was larger; adding it means increasing the exponent from $i-1$ to $i$. Since $p^{i+1}\asymp q^2>Ck$, no higher power contributes to $W_x$. The cofactors of any affected divisor below $x$ are at most $x/p^i<p$ for all large $k$. They are below $q$ and their missing harmonic mass is $o(1)$. Hence the magnitude of the change in $W_x$, divided by the one-digit cost $p-1$, is
\[
 (1+o(1))\frac{\beta_i\log q}{p^{i+1}}
       =(1+o(1))\lambda_q(t_i/p)^{i+1}.
\]
The estimate is uniform for $k\le x\le Ck$ and $p$ in such a compact range. Indeed the retained cofactor harmonic sum is $\log(x/p^i)+O(1)= (\beta_i+o(1))\log q$ and the total cutoff subtraction is $O(p^{-i})$. Distinct altered powers in any of our exchanges have product exceeding $Ck$: their exponents in $q$ are $2i/(i+1)\ge4/3$. Their gains and losses therefore add exactly before estimating.

Suppose $j\ge3$ and the formula is known through $j-1$. If $w_j\le(1-\varepsilon)t_j$ infinitely often for some $0<\varepsilon<1$, all primes in $[(1-3\varepsilon/4)t_j,(1-\varepsilon/2)t_j]$ can be raised from $j-1$ to $j$ since $w_{j-1}\gg t_j$. Their per-cost gains exceed $(1+\eta)\lambda_q$ for some fixed $\eta>0$. Their combined cost $B$ is of order $t_j^2/\log q$ by the PNT. Pay for them by trimming $(j-1)$st powers of primes just below $t_{j-1}$, in $[(1-2s)t_{j-1},(1-s)t_{j-1}]$ with fixed $s>0$ sufficiently small. Those primes have exactly $j-1$ copies and their per-cost losses are below $(1+\eta/2)\lambda_q$. The interval has available cost $\gg_s t_{j-1}^2/\log q\gg B$. Stop at the first total reaching $B$; the overshoot is $O(t_{j-1})=o(B)$ since $4/(j+1)>2/j$. This strictly improves $W_x$ uniformly, a contradiction.

If $w_j\ge(1+\varepsilon)t_j$, trim the $j$th and all higher powers in $[(1+\varepsilon/2)t_j,(1+3\varepsilon/4)t_j]$. Counting just one digit per prime frees $B\asymp t_j^2/\log q$, and the losses per one-digit cost are below $(1-\eta)\lambda_q$. Spend up to $B$ adding $(j-1)$st powers in $[(1+s)t_{j-1},(1+2s)t_{j-1}]$. Their previous exponent is exactly $j-2$ by the induction hypothesis, including $w_1\sim q$ when $j=3$. For small fixed $s$ each gain per cost exceeds $(1-\eta/2)\lambda_q$. Leave at most $O(t_{j-1})$ unspent. Again the score strictly improves. The changed periods need not be compressed: they are feasible competitors for the unrestricted period budget. The two contradictions complete the induction.
\end{proof}
